% --- Tabel Deskripsi UC-03 ---
\begin{table}[H]
    \centering
    \caption{Deskripsi \textit{Use Case} Melakukan \textit{Peer Assessment} (UC-03)}
    \label{tbl:desc_uc03}
    \begin{tabular}{|p{0.20\textwidth}|p{0.70\textwidth}|}
        \hline
        \textbf{Nama \textit{Use Case}} & \textbf{Melakukan \textit{Peer Assessment}} \\
        \hline
        \textbf{Aktor} & pelajar \\
        \hline
        \textbf{FR Terkait} & FR-03, FR-09 \\
        \hline
        \textbf{Deskripsi} & pelajar mengevaluasi pekerjaan rekan sejawat menggunakan rubrik terstruktur. Aktivitas ini dipicu oleh sistem untuk memecahkan isolasi kompetensi. \\
        \hline
        \textbf{Kondisi sebelum} & Sistem mendeteksi stagnasi pembelajaran; Algoritma \textit{Re-ranking} telah menemukan pasangan heterogen. \\
        \hline
        \textbf{Kondisi sesudah} & Skor Elo \textit{Assessee} diperbarui (berdasarkan nilai tugas); Skor Elo \textit{Assessor} diperbarui (berdasarkan kualitas reviu). \\
        \hline
        \textbf{Main Flow} & \begin{enumerate}[noitemsep,topsep=0pt,leftmargin=*]
            \item pelajar menerima notifikasi tugas \textit{peer-assessment}.
            \item Sistem menampilkan jawaban rekan \textbf{tanpa identitas (anonim)} beserta kolom penilaian.
            \item pelajar mengisi penilaian.
            \item Sistem memproses hasil untuk memperbarui Elo rating kedua belah pihak.
        \end{enumerate} \\
        \hline
    \end{tabular}
\end{table}